\documentclass[11pt,notitlepage]{report}

\usepackage[T1]{fontenc}
\usepackage{lmodern}
\usepackage[margin=1in]{geometry}
\usepackage{amsmath}
\usepackage{amssymb}
\usepackage{amsthm}
\usepackage{microtype}
\usepackage[colorlinks=true, allcolors=blue]{hyperref}

\theoremstyle{definition}
\newtheorem{definition}{Definition}
\newtheorem{algorithm}{Algorithm}

\theoremstyle{plain}
\newtheorem{lemma}{Lemma}

\setlength{\parindent}{0pt}
\setlength{\parskip}{0.5em}

\title{DQN notes}
\author{}
\date{}

\begin{document}

\maketitle

\tableofcontents

\chapter{HuggingFace Q-Learning}\label{chap:hface}
\url{https://huggingface.co/learn/deep-rl-course/en/unit2/what-is-rl}

\section{Tons of definitions}

\begin{definition}[Policy]\label{def:policy}
    Usually notated as $\pi$.
    Just a function from $States \to Actions$.

    They like to call the optimal policy $\pi^*$
\end{definition}
\begin{definition}[Policy-based method]\label{def:policy-based}
    Try to find $\pi^*$ directly.
    In other words your model will input $State$ and output $Action$.
    And you will train against some loss to get a better model.
\end{definition}
\begin{definition}[Value-based method]\label{def:value-based}
    Still trying to find $\pi^*$, but add in one intermediate step.
    Build a function that maps $States \to \mathbb{R}$. Where target is "Value".
    If your value function is optimal then you can easily get optimal policy (search neighboring states
    for highest value).
    But you can do other things based on the value too. Maybe you do a beam-search down to a certain depth.
    Maybe you set some "exploration vs exploitation" params. Etc. etc.
\end{definition}
There's some dependency between value-based and policy-based. Not inherent in the definition but pretty
common in actual practice. i.e. most of the time you look for value assuming a given policy.
(It's too hard to ascribe a direct "value" to a state. You need to know some stuff about neighbors.
But to do that you need a policy.)
\begin{definition}[Value Function]\label{def:value-function}
    Given a policy $\pi$, we make a value function.
    \[v_\pi: States \to \mathbb{R}.\]
    This particular value function uses "expected discounted return". I get the impression that this is
    a very popular way to convert policy to value.
    \[v_\pi(s) = \mathbb{E}\left[R_{t+1} + \gamma R_{t+2} + \gamma^2 R_{t+3} + \ldots \vert S_t = s\right]\]
    This is creating a value by starting in state $S_t$ and then following the policy to get to future states.
    $R_x$ is the reward you get out of state $S_x$.

    What the heck is $R_x$? They call it the "immediate reward" from the state transition $S_t \overset{action}{\to} S_{t+1}$.
    I'm guessing we'll talk how to pick "immediate reward" in the future.
    The way that makes the most sense in my head is for "reward" to only come from a terminal state, but
    I'm sure that for bigger problems that would be troublesome (ex: in binary success/fail, if your model
    never sees a success it will be REALLY HARD to learn anything).
    
    $0 \leq \gamma \leq 1$ is the "discount". My guess is that $\gamma$ is just a training param.
    Your uninitialized model will have some pretty crazy values attached to states and $\gamma$ decides
    how hard you backprop.
\end{definition}

So I guess all this setup is to motivate why do we optimize a value function instead of policy directly.
And the answer is that it helps for us to have a bit more control over the final policy.
Ex: There's a popular method called "Epsilon-Greedy" where your policy sometimes picks the best
value and sometimes doesn't (explore/exploit).

\begin{definition}[State Value Function]
    \[V_\pi(s) = E_\pi\left[G_t\vert S_t = s\right]\]
    \begin{itemize}
    \item $V_\pi(s) =$ value of state $s$ following policy $\pi$.
    \item $E_\pi\left[\ldots\right] =$ expected return following policy $\pi$.
    \item $G_t =$ They called this $v_\pi(s)$ earlier. I guess they just really
    love introducing more terms for the same thing. Anyway its the 
    $\mathbb{E}\left[R_{t+1} + \gamma R_{t+2} + \gamma^2 R_{t+3} + \ldots\right]$ you saw before.
    \item $S_t = s$ means that I'm starting in state $s$.
    \end{itemize}
\end{definition}
\begin{definition}[Action Value Function]\label{def:action-value-function}
    \[Q_\pi(s,a) = E_\pi\left[G_t\vert S_t = s, A_t = a\right]\]
    Very similar to state value but now the action at time $t$ is a param.
\end{definition}

In the website they introduce these as if they're really different.
I wonder how different they are really? $A_t$ is determined based on the policy,
and the policy is an argument of $V_\pi(s)$. I feel like one is just a weighted average
of another (assuming your policy is stochastic).

Oh I guess that's the next slide.
\begin{lemma}[Bellman Equation]
    \[V_\pi(s) = E_\pi\left[R_{t+1} + \gamma \times V_\pi\left(S_{t+1}\right)\vert S_t = s\right]\]
\end{lemma}
\begin{proof}
Ummm... Duh?
I think all that's missing is some explanation.
They don't specify what action is taken for $t+1$, but I think it's clear from context
that you get it from the policy $\pi$.
\end{proof}

\section{How to Learn?}
\begin{definition}[Monte Carlo Learning]\label{monte-carlo-learning}
    \[V(S_t) \leftarrow V(S_t) + \alpha\left[G_t - V(S_t)\right]\]
    After you finish an "episode" (whatever that is) you can update your value function.
    Just make an adjustment according to the diff between the observed return from
    the episode and the estimated return.

    And to be clear, you apply the adjustment for each step $t$ that you took in the episode.
\end{definition}
\begin{definition}[Temporal Difference Learning]\label{temporal-difference-learning}
    Don't wait until the end of the episode. You can learn at each step.
    \[V(S_t) \leftarrow V(S_t) + \alpha\left[R_{t+1} + \gamma V(S_{t+1}) - V(S_t)\right]\]
    So in the Monte Carlo above we've finished the episode so we know the exact reward $G_t$.
    Here we're still in the middle of the episode. But $R_{t+1}$ is a real signal. And $V(S_{t+1})$
    is just our best guess so far.
\end{definition}

The website provides no discussion about which method is better in which cases.
My first instinct is that Monte Carlo is really bad when your problem is big and your
policy is not well trained. If you're not going to see enough of the different terminal states
to get an idea of the shape of the problem then you're going to fail pretty badly.
But the good news in Monte Carlo is that it's easier to judge reward when you're already
at a terminal state.

I think one downside of the Temporal Difference is... who's picking that reward function?
I hope you pick correctly because the model is going to be learning entirely off of that.
If you pick a bad reward function then you'll learn garbage.
The good news though is that you can solve problems this way that you would never get
by Monte Carlo. If you model is too bad to actually "win" then it can at least get
some intermediate signal.

\section{Q-Learning}
I think this is where I get to actually learn something?
With the above context it's pretty easy to define.
You're optimizing an Action-Value function \ref{def:action-value-function}.
So you have a function $States \times Actions \to \mathbb{R}$ representing value.
And you train it using the Temporal-Difference style laid out above.
(That means you need a reward for each transition).

One common policy is called "Epsilon-Greedy".
At each step, choose the action with the highest $Q$ value with probability
$1-\epsilon$, or a random action with probability $\epsilon$.

And you'll have to cleverly vary $\epsilon$ over the course of training,
don't ask me the specifics.

\begin{definition}[On-Policy and Off-Policy]
    On-Policy = Same policy used for acting and updating
    Off-Policy = Different policy used for acting and updating
\end{definition}

Relevant here because of $\epsilon$-Greedy. The acting policy includes the
$\epsilon$ but the training policy does not.

\begin{algorithm}[Q-Learning]
\begin{itemize}
\item In state $s$, choose action $a$ (best or random).
\item Observe reward $r$ and next state $s'$.
\item Move $Q(s,a)$ toward $r + \gamma \max_{a'} Q(s',a')$ (just $r$ if terminal).
\item Repeat; restart when the episode ends.
\end{itemize}
\end{algorithm}

\section{Deep Q-Learning}

So the previous chapter explictely write a table of size $States \times Actions$
and then manually updated it (I don't include the worked example in these notes).
Which is great for games of size similar to Tic-Tac-Toe and otherwise pretty useless.

This chapter is going to talk about putting a NN in there.

They suggest that input = State, and output = Actions. Note the plural ActionSSSS.
In otherwords, the input = an encoding of the state into whatever size vector.
The output = a vector the size of all possible actions. Where the meaning is the Q-value
of that (state, action) pair.

The specific encoding I think could get interesting but we're going to ignore it for now.

They next talk about the NN structure. They're playing some video game by looking at the
screen, and they talk about convolutional layers. We're not going to do that. Skip this
bit.

\begin{definition}[Q-Target]\label{def:q-target}
    \[y_j = r_j + \gamma \text{max}_{a'}\hat{Q}\left(\phi_{j+1}, a'; \theta^-\right)\]
    Ahahaha lol what.

    I think they're saying that we would LOVE if the action NN output corresponding to
    action $a'$ corresponds to immediate reward plus discounted NN estimate of greedy
    next step.

    Unclear what $theta^-$ is. I know $theta$ is just the NN parameters.

    Oh yeah they rewrote the same thing in different notation.
    \[R_{t+1} + \gamma \text{max}_{a} Q\left(S_{t+1}, a\right)\]
    So the output of $\text{NN}(s)$ for a specific action should be the immediate
    reward gained by that action plus the NN's own prediction of greedy value from
    that state forward. (which is an estimate of the actual value from that state forward).
\end{definition}
\begin{definition}[Q-loss]\label{def:q-loss}
    \[y_j - Q\left(\phi_j, a_j; \theta\right)\]
    I'm not going to re-write it in the other notation. It's clear.
    Diff between real value and target value is loss.
\end{definition}

\section{DQN Algorithm}
Next comes the actual training algorithm:
\begin{algorithm}
It's better to train on batches, for reasons that I REALLY SHOULD look at in more detail
(IIUC, more stable gradients? Also maybe it's computationally more efficient? Or maybe not.)
So you take some number of samples and store a "replay memory", just keeping track of all
the stuff you need to know to do a train step. Then once you have enough you can do a train step.

They say (without explaining) that you should randomly sample a smaller batch from your replay
memory. I wonder why??

They also explain what the heck is $Q$ and $\hat{Q}$ above, which I appreciate.
See in \ref{def:q-target} that the loss itself depends on our estimate of value.
Well, we would not like to deal with recursive nonsense. Instead we have
both $\hat{Q}$ and $Q$.

Set $\hat{Q} = Q$ at the beginning (i.e. $\theta^- = \theta$). Train steps will
pull $Q$ away from $\hat{Q}$. Every $C$ steps reset $\hat{Q} = Q$.
\end{algorithm}

Now for some explanations!
\begin{itemize}
\item Why "replay memory". I didn't totally understand the website. One convincing reason is to have a bunch of different gamestates in each train step (more stable gradients). Otherwise, for example, if your NN is playing level 2 it will get all its gradients from level 2 and forget level 1. They also say that the NN can learn from an experience more than once. But not clear that's a positive??
\item Why $\hat{Q}$? I think we touched that above. The stated reason is to reduce oscillation. Without $\hat{Q}$ it's like the training is chasing itself. This way we chase a stable point for a while, then chase the next stable point for a while.
\item Double DQN. Okay this is really not adequately explained in the website. They're talking about having a DQN network and a target network. And they suggest decouple to prevent an overweighted Q from messing up the whole training (since we assume next step is greedy)
Is this the $\hat{Q}$ and $Q$ above? Because those aren't really decoupled... I wonder what this item is talking about.
Update: I asked chatGPT and it says this is just the $\hat{Q}$ and $Q$ above. It's supposed to reduce biases from an overestimated node screwing up all its neighbors.
I don't really agree yet ($\hat{Q}$ and $Q$ are pretty similar) but I see the intention.
\end{itemize}

Okay I think we can build our first model now?

\chapter{PyTorch DQN Tutorial}\label{chap:pytorch}

My notes on this file:
\url{https://docs.pytorch.org/tutorials/intermediate/reinforcement_q_learning.html}

\section{Two classes:}

\subsection{Transition}

A tuple representing a single transition in the environment.
Supposedly it maps
\((\text{state}, \text{action}) \rightarrow (\text{next\_state}, \text{reward})\).

\subsection{ReplayMemory}

So apparently we want to hold some amount of the most recent
\textbf{Transition} tuples.
I think it will be made clear later what this is all about. Presumably something
about the training step(s).

\section{DQN Algorithm}

\begin{definition}[Return]\label{def:return}
The \textbf{return} at time \(t_0\) is the discounted cumulative reward
\[
    R_{t_0} = \sum_{t=t_0}^{\infty} \gamma^{t-t_0} r_t.
\]
\end{definition}

Okay what are all of these things?

\begin{itemize}
    \item \(0 < \gamma < 1\). The discount on the reward. Partially to make
          the sum converge and partially to weight near rewards higher than
          far (I guess presuming your model is bad at predicting too far in
          the future)
    \item \(r_t =\) presumably the reward at transition \(t\).
\end{itemize}

\subsection{Q-learning}

Imagine we had a function
\[Q^*: State \times Action \to \mathbb{R}\]
where the value is equal to the \hyperref[def:return]{return}.

Well if we had this we could easily optimize towards this return.
\[\pi^*(s) = \text{argmax}_{a}Q^*(s,a)\]
Here $\pi^*$ is our "policy function". Given a state it picks an action.
And because we assumed $Q^*$ gives the exact return, then this $\pi^*$ is the
optimal function for that particular return function.

Well now stop assuming that we know $Q^*$ and instead train your NN to approximate it.

Apparently this thing is the "Bellman Equation". Let's puzzle out what they're talking about.
\[Q^\pi(s,a) = R(s,a) + \gamma Q^\pi(s', \pi(s'))\]
\begin{itemize}
\item $Q^\pi(s,a) =$ Expected return if we take action $a$ then use the policy $\pi$ thereafter
\item $R(s,a) =$ reward when taking action $a$ from state $s$.
\item $\gamma =$ same as above. Whatever your future discount factor.
\item $Q^\pi(s', \pi(s')) =$ hopefully clear. $s'$ is the next state after $(s,a)$.
\end{itemize}
Okay, that's not so complicated. The value of a given action is equal to its reward plus the expected
future reward from the future state(s).

The website then defines
\[\delta = Q(s,a) - \left(r+\gamma \text{max}_a' Q(s',a)\right)\]
and says we're going to minimize $\delta$. This makes no sense to me? I thought by Bellman's
equation up there we know that this $\delta$ is just zero? But maybe it's just some confusion
over what all of these $s'$, $\text{max}'$ really mean.

Okay I asked the AI for clarification. It said to clarify $Q$ above as $Q_\theta$.
\[Q_\theta = \text{neural net's prediction of } Q^*\]
So the "Bellman Equation" which makes perfect sense for a policy does not automatically apply
to $Q_\theta$.

It's a bit easier to start at the end and work backwards. In some terminal state/action $s, a$:
\[Q*(s,a) = R(s,a)\]
\[Q_\theta(s,a) = \text{depends on }\theta.\text{ But we should train it to equal } R(s,a).\]
And non-terminal states make sense too. It just depends on how $Q_\theta$ estimates the
various adjacent states.

\subsection{loss function}
They say they want to use "Huber Loss", which is essentially $\left|error\right|^2$ for small-ish values
and $\left|error\right|$ for big-ish values. I don't really care about that at this time.
We can debate different loss functions after I understand a bit better the theory.

\section{Implementation Detail}
\subsection{NN Dimensions}
Okay... so we are training a NN trying to predict $Q^*(s,a)$. How do we set up the training?
Ignore some of the NN details right now. Just say it's standard feed-forward MLP with relu.
What's the input and output?

The website says "Our model will be a feed forward neural network that takes in the difference between the current and previous screen patches"
What does that mean? I understand "feed forward NN" perfectly well. I don't understand the input.
They same the input dimension = "\texttt{n\_observations}" which... what? In their dummy problem I thought the input would
be some state of the underlying game. Where the cart is in relation to the stick or whatever.

Okay I asked Codex to explain it to me. "\texttt{n\_observations}" means the number of input features.
In this problem it's 4. Cart position, cart velocity, pole angle, pole angular velocity.

I believe the output dimension = number of possible actions
And the meaning of the output value = expected reward from this action.
For the dummy problem on the website it's simple because there are two actions and they are always available.

So after reading some this makes sense now.
Input layer = state of the game, one node per state item
Output layer = expected reward per action, one node per possible action
Middle layers = some simple MLP feed-forward RELU

\subsection{How to Train?}
I think this website I'm following is just not very good.
They introduce a "\text{target\_net}" somewhere late into it without explanation of what it's doing.
Now we have two NN. A policy NN and a target NN. I've heard of this before in other places so it
kind of makes sense but I'm disappointed in the lack of explanation. 

\chapter{Hugging Face Deep RL Course}\label{chap:huggingface}

\section{Q-Learning (Unit 2)}
\url{https://huggingface.co/learn/deep-rl-course/en/unit2/introduction}

% Add your Unit 2 notes here.

\section{Deep Q-Learning (Unit 3)}
\url{https://huggingface.co/learn/deep-rl-course/en/unit3/introduction}

% Add your Unit 3 notes here.

\end{document}
